\documentclass[11pt,a4paper]{article}
\usepackage[utf8]{inputenc}
\usepackage[T1]{fontenc}
\usepackage[spanish,es-tabla]{babel}
\usepackage[margin=2.4cm]{geometry}
\usepackage{lmodern,microtype}
\usepackage{array,tabularx,longtable,booktabs}
\usepackage{enumitem,xcolor,listings}
\usepackage{hyperref,xurl}
\hypersetup{colorlinks=true,linkcolor=blue!45!black,urlcolor=blue!60!black}
\setlength{\parindent}{0pt}
\setlength{\parskip}{0.6em}
\setlength{\emergencystretch}{3em}
\setlist[enumerate]{itemsep=0.45em,topsep=0.4em}
\newcolumntype{Y}{>{\raggedright\arraybackslash}X}
\newcommand{\campo}[1]{\nolinkurl{#1}}
\newcommand{\alerta}[1]{\begin{quote}\small\textbf{Importante:} #1\end{quote}}
\lstset{basicstyle=\ttfamily\footnotesize,columns=fullflexible,keepspaces=true,
 breaklines=true,showstringspaces=false,frame=single,rulecolor=\color{black!20}}
\title{Tutorial: web de recados con Microsoft y Azure}
\author{Next.js y React en el frontend; Python en el backend}
\date{18 de octubre de 2026}
\begin{document}
\maketitle
\tableofcontents
\newpage

\section{Tutorial}

Esta aplicación permite iniciar sesión con Microsoft y crear, consultar o borrar recados privados. Las notas se conservan en Azure Cosmos DB cuando se actualiza la web. El frontend utiliza Next.js y React; el backend utiliza Python con Flask; Docker empaqueta la aplicación.

Next.js genera archivos estáticos durante la construcción de Docker. El contenedor final ejecuta Python, que sirve esos archivos y atiende la API. Node.js participa solamente en la construcción, sin actuar como backend de la web publicada.

\subsection{Generación automática del workflow}

\textbf{El ZIP actualizado no contiene la carpeta \campo{.github/workflows}.} Primero se sube el código a GitHub y después se conecta el repositorio desde el \textbf{Centro de implementación de Azure App Service}. Azure genera el workflow y lo agrega al repositorio al guardar la integración.

No hay que crear un YAML desde GitHub ni copiar el workflow anterior. La ruta que aparecerá es \campo{.github/workflows/}: lleva punto inicial y la palabra \campo{workflows} en plural. El nombre del archivo lo determina Azure.

\alerta{Subir el código no crea por sí solo los recursos, las credenciales ni la conexión con Azure. La automatización comienza después de completar la configuración inicial.}

\subsection{Publicación de la aplicación}

\begin{center}
\setlength{\fboxsep}{8pt}
\fbox{\shortstack{GitHub\\Código fuente y Dockerfile}}\\[0.5em]
$\downarrow$\\[0.5em]
\fbox{\shortstack{GitHub Actions\\Workflow generado por Azure}}\\[0.5em]
$\downarrow$\\[0.5em]
\fbox{\shortstack{Docker\\Construcción de Next.js y Python}}\\[0.5em]
$\downarrow$\\[0.5em]
\fbox{\shortstack{Azure Container Registry\\Almacenamiento de la imagen}}\\[0.5em]
$\downarrow$\\[0.5em]
\fbox{\shortstack{Azure App Service\\Ejecución de la web con HTTPS}}
\end{center}

GitHub almacena el código y ejecuta la automatización. Azure Container Registry (ACR) guarda las imágenes Docker, pero no ejecuta la web. App Service descarga la imagen elegida y ejecuta el contenedor. Cosmos DB guarda las notas fuera del contenedor.

\subsection{Funcionamiento de la web}

\begin{enumerate}
\item El navegador solicita la página por HTTPS y Python entrega el frontend compilado.
\item Al pulsar \textbf{Continuar con Microsoft}, el backend inicia la autenticación.
\item Microsoft devuelve al usuario a \campo{/auth/callback}. Python valida la respuesta y establece la sesión.
\item React solicita las notas a la API; Python consulta Cosmos utilizando la identidad de esa sesión.
\item Al crear o borrar un recado, Python modifica solamente los datos de su propietario y React actualiza la pantalla.
\end{enumerate}

Las claves se configuran en App Service. No se guardan en GitHub, en el JavaScript del navegador ni dentro de la imagen.

\section{Requisitos y nombres del ejemplo}

\subsection{Cuentas y programas}

\begin{itemize}
\item Cuenta GitHub con permiso de administración del repositorio y GitHub Desktop instalado.
\item Suscripción Azure activa con permiso para crear los recursos del tutorial.
\item Permiso para registrar aplicaciones de Microsoft Entra, o ayuda del administrador del directorio.
\item ZIP \campo{Recados-Nextjs-Python-Azure-Automatico.zip} de esta entrega.
\end{itemize}

No hace falta instalar Docker, Python o Node.js en el computador para el recorrido principal. La primera imagen se construye con ACR Tasks desde Cloud Shell; las siguientes se construyen en GitHub Actions.

\subsection{Cuenta de estudiante y costos}

Si corresponde, activar \href{https://azure.microsoft.com/free/students/}{Azure for Students}, completar la verificación y comprobar la suscripción en \href{https://portal.azure.com/}{Azure Portal}.

Revisar regiones disponibles, cuotas y créditos antes de crear recursos. No se exige Chile Central: elegir una región admitida por la suscripción. Usar un plan compatible con contenedores Linux, por ejemplo B1 si está disponible; no asumir que F1 permite esta configuración. ACR ofrece niveles como Basic, Standard y Premium, no un plan genérico gratuito. Los recursos y construcciones pueden consumir créditos o generar cargos.

\subsection{Nombres que deben anotarse}

Estos son ejemplos; reemplazar de forma consistente los nombres no disponibles:

\begin{tabularx}{\textwidth}{@{}p{0.40\textwidth}Y@{}}
\toprule
\textbf{Elemento} & \textbf{Ejemplo}\\
\midrule
Repositorio GitHub & \campo{Tutorial-Web-II}\\
Grupo de recursos & \campo{rg-recados-web}\\
Registro ACR & \campo{registrorecados2026abc}\\
Repositorio de imágenes ACR & \campo{recados}\\
Cuenta Cosmos DB & \campo{cosmosrecados2026abc}\\
Base y contenedor de datos & \campo{RecadosDB} y \campo{recados}\\
Clave de partición & \campo{/owner}\\
Web App & \campo{web-recados-2026abc}\\
Registro de login Microsoft & \campo{recados-login}\\
Rama de publicación & \campo{main}\\
\bottomrule
\end{tabularx}

La URL pública se copia desde App Service. Puede tener sufijos adicionales. El repositorio de código de GitHub y el repositorio de imágenes de ACR son elementos distintos.

\section{Preparar y subir el proyecto a GitHub}

\subsection{Descomprimir el ZIP}

\begin{enumerate}
\item Descomprimir el ZIP y abrir la carpeta que contiene directamente \campo{Dockerfile}, \campo{server.py} y \campo{requirements.txt}.
\item Comprobar que existe \campo{frontend} con \campo{package.json}, \campo{package-lock.json} y \campo{app}.
\item Confirmar que no hay un workflow de despliegue. Azure lo agregará después.
\end{enumerate}

\begin{lstlisting}
Dockerfile
server.py
store.py
requirements.txt
.gitignore
.dockerignore
README.md
frontend/
  package.json
  package-lock.json
  next.config.mjs
  app/
    layout.jsx
    page.jsx
    globals.css
static/
templates/
tests/
docs/
\end{lstlisting}

\alerta{El Dockerfile debe quedar en la raíz del repositorio. No subir la carpeta envolvente de modo que todo quede dentro de otra subcarpeta. Subir los archivos descomprimidos, no solamente el ZIP.}

\subsection{Crear el repositorio}

\begin{enumerate}
\item Entrar en \href{https://github.com/}{GitHub} y seleccionar \textbf{New repository}.
\item Elegir propietario y nombre, por ejemplo \campo{Tutorial-Web-II}.
\item Elegir \textbf{Public} para seguir el ejemplo de construcción inicial desde una URL pública. La sección de construcción también incluye una alternativa para repositorios privados.
\item No agregar README, licencia ni \campo{.gitignore} desde el asistente: el proyecto ya contiene los archivos necesarios.
\item Pulsar \textbf{Create repository} y copiar su URL real.
\end{enumerate}

Si se reutiliza el repositorio de la calculadora, clonar ese repositorio, sustituir los archivos y retirar su antiguo despliegue. No borrar los recursos de Azure por cambiar el código.

\subsection{Subir los archivos con GitHub Desktop}

\begin{enumerate}
\item Instalar \href{https://desktop.github.com/}{GitHub Desktop} e iniciar sesión.
\item Abrir \textbf{File $\rightarrow$ Clone repository}, elegir el repositorio o pegar su URL, seleccionar una carpeta local y pulsar \textbf{Clone}.
\item Abrir \textbf{Repository $\rightarrow$ Show in Explorer}.
\item Copiar el contenido del proyecto descomprimido dentro de esa carpeta. Conservar la carpeta \campo{.git} administrada por GitHub Desktop.
\item Revisar \textbf{Changes}: deben aparecer el Dockerfile, el backend y las carpetas de la aplicación.
\item Comprobar que la rama se llame \campo{main}; si no, renombrarla desde las opciones de rama antes de conectar Azure.
\item Escribir un resumen y pulsar \textbf{Commit to main}.
\item Pulsar \textbf{Push origin} o \textbf{Publish branch}, según corresponda.
\item Abrir GitHub en el navegador y comprobar que el Dockerfile está en la primera lista de archivos.
\end{enumerate}

Todavía no habrá ejecuciones del proyecto en Actions porque no existe un workflow. No subir \campo{node_modules}, \campo{.next}, \campo{frontend/out}, entornos virtuales, archivos \campo{.env} ni credenciales; el \campo{.gitignore} excluye los archivos habituales de desarrollo.

\subsection{Si se había usado el ZIP anterior}

Retirar del repositorio el antiguo \campo{.github/workflows/azure.yml}, confirmar y enviar ese cambio antes de conectar Azure. Copiar un ZIP nuevo sobre una carpeta existente no elimina los archivos viejos. Revisar otros workflows y retirar solamente los despliegues que actualicen esta misma web; conservar las automatizaciones ajenas a ella.

\section{Crear el grupo de recursos y ACR}

\subsection{Crear el grupo de recursos}
\begin{enumerate}
\item En Azure Portal, abrir \textbf{Grupos de recursos} y pulsar \textbf{Crear}.
\item Seleccionar la suscripción y escribir \campo{rg-recados-web}, o el nombre elegido.
\item Elegir una región admitida y confirmar.
\end{enumerate}

\subsection{Crear el registro de imágenes}
\begin{enumerate}
\item Buscar \textbf{Container registries / Registros de contenedores} y pulsar \textbf{Crear}.
\item Seleccionar suscripción, grupo y región; introducir un nombre único, como \campo{registrorecados2026abc}.
\item Elegir un nivel adecuado, por ejemplo \textbf{Basic}, revisar el precio y crear el recurso.
\item Mantener una conectividad que permita construir y descargar imágenes desde los servicios utilizados. La imagen privada seguirá requiriendo autenticación.
\item Copiar el \textbf{Login server / Servidor de inicio de sesión} completo. No deducirlo solamente del nombre del registro.
\end{enumerate}

\subsection{Autenticación de la plantilla para contenedores}

La ruta de referencia es la integración de contenedores que genera Azure desde el portal. Según la plantilla disponible, puede utilizar credenciales de ACR y un perfil de publicación, o una identidad/OIDC cuando esa variante se ofrezca.

Para la plantilla basada en credenciales, abrir ACR $\rightarrow$ \textbf{Access keys / Claves de acceso}. Si el asistente exige usuario y contraseña, habilitar \textbf{Admin user / Usuario administrador}. Usar esos datos exclusivamente en los campos de autenticación de Azure que los soliciten; no agregarlos al código.

Si se elige identidad/OIDC, seguir los permisos de la sección de identidades. No combinar las referencias de secretos de un workflow OIDC anterior con una plantilla nueva basada en perfil de publicación.

\section{Construir la primera imagen Docker}

Esta construcción se realiza una vez para disponer de una imagen real al crear App Service. Las actualizaciones posteriores las construirá el workflow generado por Azure. Las cuotas y permisos de ACR Tasks deben permitir la operación.

\subsection{Desde un repositorio público}
\begin{enumerate}
\item En Azure Portal, abrir \textbf{Cloud Shell}, elegir \textbf{Bash} y completar la configuración inicial si se solicita.
\item Seleccionar la suscripción y ejecutar la construcción, reemplazando todos los datos de ejemplo por los reales.
\end{enumerate}

\begin{lstlisting}
az account set --subscription "TU_ID_DE_SUSCRIPCION"

az acr build \
  --registry "TU_NOMBRE_DE_ACR" \
  --image "recados:inicial" \
  --file "Dockerfile" \
  "https://github.com/TU_USUARIO/TU_REPOSITORIO.git#main"
\end{lstlisting}

\begin{enumerate}[resume]
\item Esperar a que termine la construcción. El Dockerfile instala las dependencias y compila Next.js antes de preparar Python.
\item Abrir ACR $\rightarrow$ \textbf{Repositories / Repositorios} y comprobar \campo{recados} con etiqueta \campo{inicial}.
\end{enumerate}

No pasar claves de Cosmos ni secretos Microsoft como argumentos de construcción: se necesitan al ejecutar la web, no al construirla.

\subsection{Alternativa para un repositorio privado}

Puede subirse el ZIP de esta entrega a Cloud Shell sin poner un token personal de GitHub en un comando:

\begin{enumerate}
\item En Cloud Shell, abrir \textbf{Upload/Download files $\rightarrow$ Upload} y subir el ZIP.
\item Descomprimirlo, entrar en la carpeta del Dockerfile y construir usando el contexto local.
\end{enumerate}

\begin{lstlisting}
unzip Recados-Nextjs-Python-Azure-Automatico.zip -d carga-recados
cd carga-recados/Recados-Nextjs-Python-Azure-Automatico

az acr build \
  --registry "TU_NOMBRE_DE_ACR" \
  --image "recados:inicial" \
  --file "Dockerfile" .
\end{lstlisting}

El punto final significa ``la carpeta actual''. Si la suscripción no permite ACR Tasks, la alternativa es construir con Docker desde un equipo que lo tenga, iniciar sesión en ACR con Azure CLI y publicar la misma imagen. Un error de cuotas o permisos no se soluciona agregando un workflow vacío.

\section{Crear Azure Cosmos DB}

\subsection{Crear la cuenta}
\begin{enumerate}
\item Buscar \textbf{Azure Cosmos DB}, pulsar \textbf{Crear} y seleccionar \textbf{Azure Cosmos DB for NoSQL}.
\item Elegir suscripción, grupo y una región admitida; asignar un nombre único, como \campo{cosmosrecados2026abc}.
\item Seleccionar la capacidad conforme al presupuesto. Serverless y el nivel gratuito de rendimiento aprovisionado son modalidades diferentes; revisar la elegibilidad y el costo mostrado.
\item Revisar y crear el recurso.
\end{enumerate}

\subsection{Crear la base y el contenedor}
\begin{enumerate}
\item Abrir \textbf{Data Explorer / Explorador de datos} y seleccionar \textbf{New Container}.
\item Crear la base \campo{RecadosDB} y el contenedor \campo{recados}.
\item Escribir exactamente \campo{/owner} como clave de partición.
\item Seleccionar el rendimiento aplicable y confirmar.
\end{enumerate}

\alerta{No reutilizar el contenedor \campo{sumas} con partición \campo{/tipo} de la calculadora. Los recados necesitan \campo{/owner}. Los nombres de base y contenedor pueden cambiar si se actualizan sus variables, pero la partición debe coincidir con el código.}

\subsection{Obtener el URI y la clave}
\begin{enumerate}
\item Abrir \textbf{Keys / Claves} en la cuenta Cosmos.
\item Copiar el URI y una clave de lectura y escritura para configurarlos después en App Service.
\item Mantener habilitada la autenticación por clave para esta implementación. No pegar esos valores en el repositorio ni en el tutorial.
\end{enumerate}

Las reglas de red se completarán cuando existan las direcciones de salida de App Service.

\section{Crear App Service con Docker}

\subsection{Crear la Web App}
\begin{enumerate}
\item Buscar \textbf{App Services $\rightarrow$ Crear $\rightarrow$ Aplicación web}.
\item Seleccionar suscripción y grupo; asignar un nombre único a la web.
\item En \textbf{Publicar}, elegir \textbf{Contenedor / Container}; no publicación directa de código Python.
\item Elegir \textbf{Linux}, una región admitida y un plan compatible con contenedores, como B1 si está disponible.
\item Para seguir la ruta de referencia de esta guía, usar \textbf{contenedor único clásico}. Si aparece configuración mejorada con sidecar, desactivarla para esta instalación nueva. El proyecto no necesita contenedores adicionales.
\item Elegir \textbf{Azure Container Registry}, el registro creado, la imagen \campo{recados} y etiqueta \campo{inicial}.
\item Completar la autenticación del registro. Dejar vacío el comando de inicio para conservar el del Dockerfile.
\item Revisar el precio y crear la aplicación.
\end{enumerate}

Si se reutiliza una Web App con sidecars, consultar la sección específica antes de cambiar su configuración.

\subsection{Dominio y puerto}
\begin{enumerate}
\item En \textbf{Overview / Información general}, copiar el \textbf{Default domain / Dominio predeterminado} completo y formar su URL con HTTPS.
\item Habilitar \textbf{HTTPS Only / Solo HTTPS}.
\item Agregar la configuración de aplicación \campo{WEBSITES_PORT=8080}, correspondiente al modelo clásico elegido.
\end{enumerate}

La imagen inicial puede no mostrar todavía la aplicación porque faltan las variables y credenciales de ejecución. No confundir este estado temporal con un error de construcción.

\section{Permitir la conexión de App Service con Cosmos}
\begin{enumerate}
\item En las propiedades de App Service, localizar las IP de salida y las direcciones de salida adicionales o posibles.
\item En Cosmos, abrir \textbf{Networking / Redes}.
\item Para la instalación sin red privada, usar acceso de red pública con \textbf{redes seleccionadas} y autorizar las direcciones de salida necesarias de App Service.
\item Para administrar documentos desde el portal, habilitar su acceso o agregar la IP del administrador según las opciones disponibles.
\item Guardar y esperar a que se propaguen las reglas.
\end{enumerate}

La clave no evita el firewall. Revisar las reglas si cambian las IP de salida. Una instalación con endpoints privados requiere además integración de red y DNS; esa variante no forma parte del recorrido principal.

\section{Configurar el login con cuentas Microsoft}

Este registro identifica a los visitantes. Es independiente de las credenciales del despliegue.

\begin{enumerate}
\item Abrir \textbf{Microsoft Entra ID $\rightarrow$ App registrations $\rightarrow$ New registration}.
\item Escribir \campo{recados-login} y seleccionar \textbf{Cuentas de cualquier directorio organizativo y cuentas personales de Microsoft}.
\item Elegir plataforma \textbf{Web} y la URI \campo{https://TU_DOMINIO_REAL/auth/callback}.
\item Registrar y copiar el \textbf{Application (client) ID}.
\item Abrir \textbf{Certificates \& secrets}, crear un secreto de cliente y copiar su \textbf{valor}, no su identificador. Anotar el vencimiento.
\item En \textbf{Authentication}, comprobar la URI de callback y agregar \campo{https://TU_DOMINIO_REAL/} entre las URI admitidas para el retorno después de cerrar sesión cuando Microsoft lo requiera.
\end{enumerate}

No activar flujo implícito. Python utiliza MSAL y código de autorización; no necesita leer correos ni acceder a Microsoft Graph para guardar recados. Si hay restricciones de consentimiento o de registro, debe intervenir el administrador del directorio.

\section{Configurar las variables de ejecución}

Abrir App Service $\rightarrow$ \textbf{Configuración $\rightarrow$ Variables de entorno $\rightarrow$ Configuración de aplicación}. Agregar y guardar:

\begin{longtable}{@{}p{0.35\textwidth}p{0.60\textwidth}@{}}
\toprule
\textbf{Nombre} & \textbf{Valor}\\
\midrule
\endfirsthead
\toprule
\textbf{Nombre} & \textbf{Valor}\\
\midrule
\endhead
\campo{SECRET_KEY} & Cadena aleatoria estable de al menos 32 caracteres.\\[0.45em]
\campo{APP_URL} & URL HTTPS pública exacta, sin barra final.\\[0.45em]
\campo{MICROSOFT_CLIENT_ID} & Id. de cliente de \campo{recados-login}.\\[0.45em]
\campo{MICROSOFT_CLIENT_SECRET} & Valor del secreto de ese registro.\\[0.45em]
\campo{MICROSOFT_TENANT} & \campo{common}.\\[0.45em]
\campo{COSMOS_ENDPOINT} & URI de la cuenta Cosmos.\\[0.45em]
\campo{COSMOS_KEY} & Clave real de lectura y escritura.\\[0.45em]
\campo{COSMOS_DATABASE} & \campo{RecadosDB}.\\[0.45em]
\campo{COSMOS_CONTAINER} & \campo{recados}.\\[0.45em]
\campo{WEBSITES_PORT} & \campo{8080}, para el modelo clásico.\\
\bottomrule
\end{longtable}

Para generar la clave de sesión, ejecutar en Cloud Shell:
\begin{lstlisting}
python3 -c "import secrets; print(secrets.token_hex(32))"
\end{lstlisting}
Copiar el resultado únicamente en \campo{SECRET_KEY}. Mantenerla estable entre despliegues e instancias.

\begin{enumerate}
\item Guardar las variables y aceptar el reinicio.
\item Mantener desactivada la autenticación integrada de App Service (Easy Auth): esta web gestiona su propio acceso mediante Python.
\item No sobrescribir el comando de arranque del Dockerfile.
\item Abrir \campo{https://TU_DOMINIO_REAL/health}. Debe responder \campo{status: ok}.
\end{enumerate}

No se utilizan \campo{DB_BACKEND}, \campo{SQLITE_PATH} ni \campo{TRUST_PROXY_HEADERS} de versiones anteriores. La aplicación no tiene SQLite ni base local. No crear variables \campo{NEXT_PUBLIC_} con secretos: ese prefijo puede exponerlos al navegador.

\section{Hacer que Azure genere el workflow}

\subsection{Conectar GitHub desde Azure}
\begin{enumerate}
\item En App Service, abrir \textbf{Deployment Center / Centro de implementación $\rightarrow$ Settings / Configuración}.
\item En \textbf{Source / Origen}, seleccionar \textbf{GitHub Actions} para construir la imagen desde el código.
\item Pulsar \textbf{Authorize / Autorizar}, iniciar sesión en GitHub y permitir el acceso al repositorio.
\item Elegir organización o usuario, repositorio y rama \campo{main}.
\item Elegir \textbf{Azure Container Registry}, el registro creado y la imagen \campo{recados}. Completar imagen, etiqueta y autenticación según lo que solicite el asistente. La imagen inicial ya existe.
\item En \textbf{Workflow Configuration}, pulsar \textbf{Preview file / Vista previa del archivo}.
\item Revisar los puntos de la siguiente subsección y pulsar \textbf{Save / Guardar}.
\end{enumerate}

\alerta{Seleccionar \textbf{Container registry} como único origen puede configurar un webhook que reacciona a nuevas imágenes, pero no necesariamente construye el código cuando cambia GitHub. En este tutorial se selecciona \textbf{GitHub Actions} como origen del código.}

Si la opción no aparece, comprobar que se está en la configuración de despliegue y no únicamente en la edición de imagen de la pestaña de contenedores. Las opciones pueden variar entre modelo clásico y sidecar. No sustituir la publicación de contenedor por una plantilla de código Python.

\subsection{Qué revisar en la vista previa}

El archivo debe descargar el repositorio, construir con el Dockerfile de la raíz y contexto \campo{.}, autenticarse en ACR, publicar la imagen y actualizar la Web App. Debe dispararse cuando se envían cambios a \campo{main}.

La construcción de Next.js ya está dentro del Dockerfile; la plantilla no necesita ejecutar \campo{npm run build} por separado. Tampoco debe tratar el proyecto como una web Node.js que necesite servidor Next.js en producción.

Si la vista previa solo comprime y publica código mediante un parámetro \campo{package}, sin construir Docker, corresponde a otra modalidad. Volver al asistente de contenedores y elegir la integración correcta antes de guardar.

\subsection{Credenciales que puede crear Azure}

En la plantilla de contenedor basada en credenciales, Azure agrega a GitHub los secretos del registro y el perfil de publicación. Si el portal requiere \textbf{SCM Basic Auth Publishing Credentials} para obtener o utilizar ese perfil, habilitarlo en la configuración de publicación de App Service.

Esa opción autoriza la publicación; no cambia el login Microsoft de los usuarios. Los nombres de secretos deben coincidir con las referencias del archivo generado. No inventarlos ni crear los tres \campo{AZURE_CLIENT_ID}, \campo{AZURE_TENANT_ID} y \campo{AZURE_SUBSCRIPTION_ID} del workflow anterior salvo que la plantilla real utilice OIDC y los solicite.

Si las políticas de la organización bloquean perfiles o credenciales básicas, usar una integración de contenedor que ofrezca identidad/OIDC y sus permisos. No desactivar una política de la organización para utilizar una plantilla que no admite.

\subsection{Confirmar que Azure agregó el workflow}
\begin{enumerate}
\item Abrir el repositorio en GitHub y confirmar el nuevo commit de Azure.
\item Comprobar el archivo dentro de \campo{.github/workflows/} y anotar su nombre real; no tiene por qué llamarse \campo{azure.yml} ni \campo{main.yml}.
\item Abrir \textbf{Actions} y localizar el flujo de la Web App.
\item En \textbf{Settings $\rightarrow$ Secrets and variables $\rightarrow$ Actions}, comprobar que existen los secretos referenciados. No poder leer sus valores después de guardarlos es normal.
\item En GitHub Desktop, pulsar \textbf{Fetch origin} y después \textbf{Pull origin} para descargar el commit que generó Azure.
\end{enumerate}

No recrear ni borrar el archivo generado para cada actualización. Si una protección de rama impide que Azure escriba el workflow, el administrador debe permitir la operación según las reglas del repositorio. No se resuelve creando una carpeta vacía a mano.

\section{Identidades y permisos si la integración usa OIDC}

Esta sección solo aplica si la variante de contenedor del asistente ofrece identidad/OIDC. No hay que combinarla con la ruta basada en perfil.

\begin{enumerate}
\item Seleccionar una identidad asignada por el usuario en la integración cuando se ofrezca. Con permisos suficientes, Azure puede crear la relación federada y los secretos usados por su plantilla.
\item Verificar el permiso para actualizar la Web App: habitualmente \campo{Website Contributor} en el ámbito de esa aplicación.
\item Para publicar imágenes en ACR, conceder a la identidad del despliegue \campo{AcrPush} en RBAC clásico, o \campo{Container Registry Repository Writer} si el registro utiliza ABAC.
\item Si el workflow consulta metadatos ARM del registro, conceder la lectura necesaria, por ejemplo \campo{Reader} sobre ACR.
\item Si App Service descarga mediante identidad administrada, asociarla a la web, seleccionarla en la autenticación de extracción y conceder \campo{AcrPull}, o \campo{Container Registry Repository Reader} en ABAC.
\item Si aparecen errores de token o audiencia al descargar desde ACR, revisar también la configuración de aceptación de tokens del registro en la documentación de identidad administrada.
\end{enumerate}

Publicar una imagen y descargarla son funciones diferentes, aunque una identidad pueda recibir ambas. Ninguna sustituye al registro \campo{recados-login} de los visitantes. El backend de esta entrega utiliza \campo{COSMOS_KEY} para Cosmos.

Asignar roles exige permisos de administración de acceso. Si Azure no permite crear o elegir una identidad, solicitar al administrador la identidad y los roles necesarios. No asumir que un Contributor ordinario permite administrar todas las asignaciones.

\section{Si la Web App existente usa sidecar}

El código funciona dentro de un contenedor principal con puerto 8080, pero el despliegue debe coincidir con el modelo real del servicio:

\begin{enumerate}
\item En \textbf{Centro de implementación $\rightarrow$ Contenedores}, identificar el contenedor principal y su nombre exacto.
\item Configurar el puerto de destino como \campo{8080}. En sidecar, \campo{WEBSITES_PORT} no sustituye ese ajuste.
\item Revisar que la integración generada admita ese modelo y actualice el contenedor principal correcto.
\item Mantener la autenticación y los permisos de la identidad que descarga la imagen.
\item Si el portal no ofrece un generador compatible en ese modo, no asumir que una plantilla clásica lo actualizará. Para seguir esta guía sin mantener un YAML propio, utilizar una nueva Web App de contenedor único con la integración descrita, si está disponible en la suscripción.
\end{enumerate}

No se incluye un workflow alternativo en el ZIP: Azure es quien debe generarlo. La disponibilidad de las variantes depende de la configuración y permisos del portal. No borrar ni convertir automáticamente una instalación existente.

\section{Comprobar el primer despliegue}
\begin{enumerate}
\item Abrir \textbf{GitHub $\rightarrow$ Actions} y la ejecución del workflow que agregó Azure.
\item Revisar construcción Docker, publicación de imagen y despliegue.
\item Si no comenzó una ejecución, revisar el evento y la rama del archivo. Si incluye \campo{workflow_dispatch}, usar \textbf{Run workflow}; de lo contrario, enviar un cambio pequeño a \campo{main}.
\item Abrir la URL HTTPS de App Service y comprobar la página de Recados.
\item Abrir \campo{/health} y verificar \campo{status: ok}.
\item Iniciar sesión con Microsoft, guardar una nota, cerrar sesión y entrar de nuevo. La nota debe permanecer.
\item Entrar con otra cuenta y comprobar que no ve la nota de la primera.
\item Crear y borrar una nota desde la interfaz.
\end{enumerate}

Un workflow verde confirma sus propios pasos, pero también hay que revisar el arranque del contenedor y la conexión con Cosmos. La plantilla que genere Azure puede no incluir las pruebas unitarias ni una comprobación de salud. No asumir esas verificaciones si no aparecen en su YAML.

\section{Actualizar la web posteriormente}
\begin{enumerate}
\item En GitHub Desktop, usar \textbf{Fetch origin} y \textbf{Pull origin} antes de editar.
\item Modificar los archivos, por ejemplo \campo{frontend/app/page.jsx} para la interfaz o \campo{server.py} para el backend.
\item Revisar cambios, escribir un mensaje y pulsar \textbf{Commit to main}.
\item Pulsar \textbf{Push origin}, esperar el workflow y comprobar la web.
\end{enumerate}

También puede editarse un archivo desde GitHub y confirmar el cambio en la rama configurada. Si se trabaja con otras ramas, integrar los cambios en la rama de publicación para activar el evento correspondiente.

\textbf{Guardar un archivo local no lo publica: hace falta confirmar y enviar el cambio.} No repetir ACR Tasks, recrear recursos ni volver a conectar el Centro de implementación para cada edición.

Las notas permanecen en Cosmos mientras no se eliminen sus documentos o recursos. No borrar la base para cambiar la interfaz. Para volver a una versión previa, seleccionar una imagen conservada en ACR que sea compatible con el esquema de datos.

\section{Resolver problemas habituales}

\subsection{No aparece el workflow}
Comprobar que se guardó la conexión con GitHub Actions, que Azure puede escribir en el repositorio y que se mira la rama correcta. Revisar las protecciones de rama y las políticas de Actions. Una conexión solo con ACR no equivale a construir desde GitHub.

\subsection{Falla la construcción}
Comprobar el Dockerfile en la raíz y la presencia de \campo{frontend/package-lock.json}. Revisar cuotas de ACR Tasks, acceso a imágenes base y el error concreto del log. No subir el proyecto dentro de una carpeta adicional.

\subsection{App Service no descarga la imagen}
Revisar registro, repositorio, etiqueta, red y credenciales o identidad. Poder publicar desde GitHub no implica poder descargar desde App Service. Si se utiliza identidad, revisar sus roles de extracción y la autenticación del registro.

\subsection{Errores 502, 503 o contenedor que no arranca}
Revisar logs, variables, puerto 8080 y comando de inicio. Comprobar que Next.js se construyó. Para Cosmos, revisar URI, clave, base, contenedor, partición \campo{/owner} y firewall.

\subsection{Error Microsoft AADSTS50011}
Comparar el dominio HTTPS y \campo{/auth/callback} con la URI registrada en Entra y \campo{APP_URL}. Copiar todos los sufijos del dominio. No utilizar la URL de GitHub como retorno de la web.

\subsection{Error de sesión, estado o CSRF}
Volver a entrar desde el dominio correcto, comprobar cookies y mantener una \campo{SECRET_KEY} estable. No alternar dominios durante el acceso. Recargar la página si cambió la sesión antes de crear o borrar una nota.

\subsection{Falla la autenticación del workflow}
Abrir el archivo generado y comprobar el método real: perfil, credenciales de registro u OIDC. Revisar los secretos referenciados y las credenciales que pudieron haberse revocado. No añadir secretos antiguos que el archivo actual no utiliza.

\section{Resumen de tareas}
\begin{tabularx}{\textwidth}{@{}p{0.25\textwidth}Y@{}}
\toprule
\textbf{Momento} & \textbf{Acciones}\\
\midrule
Una vez & Subir código; crear recursos; construir imagen inicial; configurar Microsoft, variables y redes; conectar GitHub desde Azure.\\[0.5em]
Cada cambio & Editar, confirmar, enviar a la rama configurada y comprobar workflow y web.\\[0.5em]
Mantenimiento & Revisar costos, respaldos, vencimiento de secretos y permisos.\\
\bottomrule
\end{tabularx}

No se ha conectado ninguna cuenta real ni desplegado recursos al preparar este documento o el ZIP.

\section{Documentación y referencias}

Las instrucciones se basan en el proyecto entregado y en estas referencias oficiales. Las pantallas y plantillas pueden variar: revisar el archivo que genere la configuración real.
\begin{itemize}
\item CI/CD de contenedores y generación desde Azure:\\
\url{https://learn.microsoft.com/en-us/azure/app-service/deploy-ci-cd-custom-container}
\item GitHub Actions y Centro de implementación:\\
\url{https://learn.microsoft.com/en-us/azure/app-service/deploy-github-actions}
\item Configuración del despliegue continuo:\\
\url{https://learn.microsoft.com/en-us/azure/app-service/deploy-continuous-deployment}
\item Construcciones con ACR Tasks:\\
\url{https://learn.microsoft.com/en-us/azure/container-registry/container-registry-tutorial-quick-task}
\item Referencia de ACR en Azure CLI:\\
\url{https://learn.microsoft.com/en-us/cli/azure/acr}
\item Puertos, variables e identidades:\\
\url{https://learn.microsoft.com/en-us/azure/app-service/configure-custom-container}
\item Configuración de sidecars:\\
\url{https://learn.microsoft.com/en-us/azure/app-service/configure-sidecar}
\item Autenticación de Microsoft con Python:\\
\url{https://learn.microsoft.com/en-us/entra/msal/python/getting-started/acquiring-tokens}
\item Exportación estática de Next.js:\\
\url{https://nextjs.org/docs/app/guides/static-exports}
\item Workflows de GitHub:\\
\url{https://docs.github.com/en/actions/concepts/workflows-and-actions/workflows}
\end{itemize}

\end{document}
